\section{Reconciliation with the construction-era model}\label{sec:upreconcile}

\subsection{The catalogue formula is reproduced}

The construction-era model sized the Qserv database in
\tabref{tab:datasetSizingOps} as
\[
  \bigl(N_{\mathrm{obj}}\,(r_{\mathrm{obj}} + r_{\mathrm{obj,extra}})
        + N_{\mathrm{src}}\,r_{\mathrm{src}}
        + N_{\mathrm{fs}}\,r_{\mathrm{fs}}\bigr) \times R ,
\]
which gives 3{,}509.56\,TB at LOY1. Multiplied by the overlap uplift of
\smQsOverlap{}, that is exactly the 2026 year-2 figure of \smQsDRoneDiskTB{}\,TB.

The two models share LDM-141 row counts and three-fold replication. The match
therefore shows that the 2026 model follows the same arithmetic; it does not
independently validate LDM-141. It also pins down three things:
\begin{itemize}
\item Neither model applies compression to the catalogue volume.
\item The 2026 model adds the measured overlap; the construction-era model
      had none.
\item The construction-era LOY1 is a full survey year. It corresponds to the
      2026 year-2 column, not the DP2 column (year 1).
\end{itemize}

\subsection{Why the 2026 demand and cost differ}\label{sec:upwhy}

\tabref{tab:smCompareVone} sets the construction-era operations demand, as
published, against the 2026 model at the same number of survey years. Both
assume all processing at the USDF. Most demand is lower in 2026, but not all
of it, and the reasons are specific.

\smpara{Raw input: \smCmpInputOne{} of the construction-era figure.}
The construction-era tables give 1{,}911\,TB a year of compressed input from
300 nights of 1000 visits, or 6.37\,GB per visit. They counted 4 bytes per pixel
and kept two images per visit as the baseline (\secref{sec:sizeinputs}). The
2026 model has 700 visits a night of 8.19\,GB logical at a compression of
0.464, or 3.80\,GB per visit. Fewer visits (0.70) times less data per visit
(0.60) gives 0.42.

\smpara{DRP compute: \smCmpComputeOne{} at one survey year, \smCmpComputeNine{} at nine, in
core-hours.}
Less input is partly offset by a higher rate: \smChPerTB{} core-hours per input
TB now, against about 23{,}500 implied by \tabref{tab:computeSizingOps}. That
rate was scaled from HSC processing plus allowances for missing pipeline steps.
The two figures are also in \emph{different units}: the construction-era
core-hour is a 2015 Intel Xeon E5-2680v3 core, and the 2026 core-hour is a
Milano core, which does more work per hour. The fall in core-hours therefore
overstates any fall in the processing work itself.

\smpara{On-floor disk: \smCmpDiskOne{}--\smCmpDiskNine{}.}
Raw data on the floor is only \smCmpRawNine{} of the construction-era figure
at nine survey years, following the smaller input. The rest of the disk is structured
differently:
\begin{itemize}
\item The construction era held every significant intermediate on disk until
      each campaign finished, kept catalogues from three releases in object
      storage, and itemised science-user space.
\item The 2026 model is dominated by live intermediates at an assumed
      \smLiveFrac{} of each campaign's intermediates, plus a
      \smOtherFrac{} allowance.
\end{itemize}
The comparison is therefore only as good as the live-intermediate fraction,
which is unmeasured.

\smpara{Tape: \smCmpTapeOne{} at one survey year, \smCmpTapeNine{} at nine.}
This is the largest difference, and it comes from policy rather than
measurement. The construction-era model archived all raw images and all data
products, and also backed up every filesystem to tape, including annual
snapshots (\secref{sec:sizeinputs}). The 2026 model archives raw data and the
retained release products only. It adds no new backup of on-floor data,
because a separate backup system already exists. If filesystem backups were
restored to scope, tape would grow substantially.

\smpara{Qserv: \smCmpQservOne{} at one survey year, rising to \smCmpQservNine{}.}
This is the one demand that is \emph{larger}, and the difference is the
overlap allowance. Both models hold three releases in Qserv, and both apply
the same catalogue formula to each release (\secref{sec:upreconcile}). The
2026 model then multiplies each release by the measured overlap uplift of
\smQsOverlap{}; the construction-era model has no such term.

The ratio is lower at one survey year for two reasons. The construction-era first year
also carries 584\,TB beyond its first release (4{,}094\,TB against 3{,}510\,TB).
The 2026 year-2 window includes the DP2 catalogue. Node counts follow storage,
because at these points both models assume the same usable data per node.

\smpara{Cost.}
Lower demand does not mean lower cost. \tabref{tab:opsSummary} put USDF
hardware at about \$57M over ten years, after deducting three quarters of the
DRP compute cost for IN2P3 and UKDF. The comparable 2026 figure, the shared-DRP forecast
(\secref{sec:upbudget}), is \smSplitTenYearHardware{} for hardware. The
differences are in the prices:
\begin{itemize}
\item Compute: about \$102 per core for the ``large Rome'' node used in the
      construction-era operations tables ($\$13{,}000$ for 128 cores), against
      \smTorinoUSD{} for a 120-core Torino node, about \$745 per core, before
      any difference in per-core speed.
\item Disk: \$70 per TB for normal filesystem disk, against \smCephPerTB{} per
      TB of usable Ceph.
\item Qserv: \$20{,}000 per node, against \smKeightsUSD{}.
\item Tape is the exception, and is cheaper: \$15.65 per TB against
      \$\smTapePerTB{}.
\end{itemize}
The construction-era model also assumed prices would fall every year (10\% for
CPU, 5\% for storage, 8\% for Qserv), where the 2026 model holds them flat. Its
own text notes that a 15\% annual storage fall would halve the operations
estimate. The longer 2026 refresh cycle (\smLifeCompute{} years against three)
works the other way, lowering cost.

% GENERATED by generate_model.py from lsst/rubin-sizing-model @ v2026.09.2. Do not edit by hand.
% Regenerate: git clone --branch v2026.09.2 https://github.com/lsst/rubin-sizing-model ; then, inside the clone, uv run generate_model.py --emit-tex <this directory>

\tiny \begin{longtable}{|p{0.20\textwidth}|l|r|r|r|r|r|r|r|r|r|}
\caption{Construction-era demand (DMTN-135 as published in 2023) against the 2026 model, at the same number of survey years n: the construction-era LOYn against the 2026 DRn, all processing at the USDF. Ratio is 2026 over 2023. \label{tab:smCompareVone}}\\
\hline
\textbf{Metric} & \textbf{Unit} & \textbf{2023 n=1} & \textbf{2026 n=1} & \textbf{ratio} & \textbf{2023 n=5} & \textbf{2026 n=5} & \textbf{ratio} & \textbf{2023 n=9} & \textbf{2026 n=9} & \textbf{ratio} \\
\hline\endfirsthead
\hline
\textbf{Metric} & \textbf{Unit} & \textbf{2023 n=1} & \textbf{2026 n=1} & \textbf{ratio} & \textbf{2023 n=5} & \textbf{2026 n=5} & \textbf{ratio} & \textbf{2023 n=9} & \textbf{2026 n=9} & \textbf{ratio} \\
\hline\endhead
Raw input, cumulative & TB & 1,911 & 798 & 0.42 & 9,556 & 3,991 & 0.42 & 17,200 & 7,184 & 0.42 \\ \hline
DRP compute & M core-hours & 45.0 & 35.5 & 0.79 & 200.0 & 177.4 & 0.89 & 370.0 & 319.3 & 0.86 \\ \hline
On-floor disk, excluding Qserv & TB & 40,698 & 30,372 & 0.75 & 190,257 & 152,557 & 0.80 & 342,575 & 276,883 & 0.81 \\ \hline
of which raw data & TB & 6,982 & 2,234 & 0.32 & 26,246 & 6,702 & 0.26 & 45,510 & 11,170 & 0.25 \\ \hline
Tape footprint & TB & 40,472 & 12,218 & 0.30 & 288,456 & 41,662 & 0.14 & 768,654 & 99,650 & 0.13 \\ \hline
Qserv storage, release window & TB & 4,094 & 4,442 & 1.08 & 31,277 & 38,609 & 1.23 & 63,206 & 77,871 & 1.23 \\ \hline
Qserv nodes & nodes & 95 & 103 & 1.08 & 364 & 447 & 1.23 & 367 & 451 & 1.23 \\ \hline
\end{longtable} \normalsize


\subsection{Superseded figures}

\begin{itemize}
\item Hosting site: NCSA $\rightarrow$ SLAC S3DF.
\item Visits per night: 1000 $\rightarrow$ \smVisits{}, over \smNights{}
      observing nights a year (\tabref{tab:upChanged}).
\item DR1 DRP compute: design-era estimate $\rightarrow$ \smDRoneNodeDays{}
      node-days for \smDRoneVisits{} visits, measured on \vthirty{}
      (\secref{sec:upcompute}).
\item Compute architecture: Xeon and Rome pricing exemplars $\rightarrow$ the
      deployed AMD EPYC Milano and Torino batch nodes and EPYC 7543P service
      nodes.
\item International split: IN2P3 50\% and UKDF 25\%, flat $\rightarrow$
      France \smFrShare{}; UK front-loaded, tapering to a \smUkShare{}
      aggregate (\secref{sec:upsplit}).
\item Hardware refresh: spread evenly $\rightarrow$ by cohort, with a
      \smLifeCompute{}-year lifetime (\secref{sec:upbudget}).
\end{itemize}

Unchanged: Qserv three-fold replication, LDM-141 catalogue counts and row
sizes, 300 observing nights, and a three-release Qserv window.

\subsection{Not updated}

The following remain as in the construction-era sections, as listed in
\secref{sec:upstatus}:
\begin{itemize}
\item the Chilean Base Data Centre and Chilean Data Access Centre (storage,
      compute and costing);
\item cloud costs;
\item a price per TB for the fast (NVMe) tier;
\item the split between normal disk and object store;
\item floating-point ratings of machines;
\item network and WAN, which were never tabulated in this document.
\end{itemize}
